% ============================================================
%  all-frames-tex.tex (Consolidated Manuscript Draft - v11)
%  A Unified Framework: Numerical Damping, Geometric Ratios,
%  Algebraic Symmetry-Breaking, and Spectral Theory
%  Author: Abdulla M. N. Al-Jabri
%  October 2026
%  License: CC-BY-4.0
% ============================================================

\documentclass[11pt,a4paper]{article}

\usepackage[utf8]{inputenc}
\usepackage[T1]{fontenc}
\usepackage{amsmath,amssymb,amsthm}
\usepackage{geometry}
\usepackage{hyperref}
\usepackage{booktabs}
\usepackage{microtype}
\usepackage{fancyhdr}
\usepackage{graphicx}
\usepackage{float}

\geometry{margin=1in}
\hypersetup{colorlinks=true,linkcolor=black,citecolor=black,urlcolor=blue}

% ---------- Footer and Header ----------
\pagestyle{fancy}
\fancyhf{}
\fancyhead[L]{\small\textit{Al-Jabri (2026)}}
\fancyhead[R]{\small\textit{The Al-Hindara Unified Framework}}
\fancyfoot[L]{\small Al-Hindara Project — Consolidated Draft}
\fancyfoot[C]{\small\thepage}
\fancyfoot[R]{\small DOI: [Pending Verification]}
\renewcommand{\headrulewidth}{0.4pt}
\renewcommand{\footrulewidth}{0.4pt}

\fancypagestyle{plain}{
  \fancyhf{}
  \fancyfoot[L]{\small Al-Hindara Project — Consolidated Draft}
  \fancyfoot[C]{\small\thepage}
  \fancyfoot[R]{\small DOI: [Pending Verification]}
  \renewcommand{\headrulewidth}{0pt}
  \renewcommand{\footrulewidth}{0.4pt}
}

\newtheorem{proposition}{Proposition}
\newtheorem{corollary}[proposition]{Corollary}
\newtheorem{observation}[proposition]{Observation}
\newtheorem{remark}{Remark}

\title{\textbf{A Unified Framework: Numerical Damping, Geometric Ratios,\\
Algebraic Symmetry-Breaking, and Spectral Theory}}

\author{
    Abdulla M.~N.~Al-Jabri\thanks{Independent Researcher, Sana'a, Yemen. ORCID: 0009-0003-3319-3822. Email: \texttt{jabri62018@gmail.com}.} \\
    \small Independent Researcher \\
    \small Sana'a, Yemen
}

\date{October 2026}

\begin{document}

\maketitle
\thispagestyle{plain}

\begin{abstract}
This paper consolidates four exploratory frameworks of the Al-Hindara project. While presented together, these paths do not constitute a single continuous chain of mathematical proof; instead, they represent four distinct exploratory studies presented within a common methodological framework. We define a numerical damping function demonstrating specific accumulation properties. We derive the geometric ratio \(\sqrt{2}-1\) and present exploratory field measurements. We establish an algebraic gauging identity and demonstrate numerically how structural asymmetry breaks symmetry-induced cancellation. Finally, we formulate a spectral problem where the kernel acts as a formal zero mode. We maintain a strict boundary between analytically proven identities, numerically computed results (independently reproduced using the supplied Python code), and empirical observations.
\end{abstract}

\vspace{1em}
\noindent\textbf{Keywords:} Sturm-Liouville theory; symmetry-breaking; numerical integration; geometric ratios; Riemann zeta function.

\tableofcontents
\vspace{1em}
\newpage

% ============================================================
\section{Introduction and Methodology}
\label{sec:intro}
% ============================================================

The Al-Hindara project explores mathematical structures across numerical, geometric, and algebraic domains. The central methodology of this paper is the rigorous separation of claims into three categories:
\begin{enumerate}
    \item \textbf{Analytically Proven:} Exact identities and algebraic relations derived from explicit definitions.
    \item \textbf{Numerically Computed and Verified:} Reproducible computations and convergence behaviors for specified parameters.
    \item \textbf{Empirically Observed:} Exploratory data presented without claims of universal physical laws.
\end{enumerate}

\textbf{Scope of Integration:} It is crucial to note that the four sections of this paper represent parallel exploratory paths. While they share the Al-Hindara context, the geometric observations (Section~\ref{sec:geometric}) neither mathematically cause nor prove the algebraic symmetry-breaking (Section~\ref{sec:algebraic}), nor do they dictate the spectral properties (Section~\ref{sec:spectral}). The unifying element is the methodology, not a cascading mathematical proof.

% ============================================================
\section{The Numerical Damping Framework}
\label{sec:numerical}
% ============================================================

\subsection{Definition and Sign Analysis}
We consider the real function \(Z(x; x_p)\) defined on \(x \ge 1\):
\begin{equation}
\label{eq:Z}
Z(x; x_p) = x^{-5} \ln(x) \sin\left(\frac{2\pi}{x}\right) e^{-x/x_p}.
\end{equation}
The parameter \(x_p\) serves as a damping parameter. The specific choice \(x_p = 21 = \lfloor \gamma_2 \rfloor\), where \(\gamma_2 \approx 21.022\) is the imaginary part of the second positive non-trivial zero of the Riemann zeta function, is a \emph{modelling choice}. No arithmetic property of \(\gamma_2\) enters the analytical proofs.

For \(x_p > 0\), the sign of \(Z(x)\) is dictated by \(\sin(2\pi/x)\). Thus, \(Z(x) < 0\) on \((1, 2)\), \(Z(2) = 0\), and \(Z(x) > 0\) on \((2, \infty)\).

\subsection{The Accumulation Function}
Define the accumulation function \(A(x) := \int_1^x Z(t; 21) \, dt\). Since \(A'(x) = Z(x)\), it follows that \(A(x)\) strictly decreases on \((1, 2)\), reaches a global minimum at \(x = 2\), and strictly increases for all \(x > 2\). (An initial hypothesis of a local maximum at \(x = 42\) is analytically false).

% ============================================================
\section{Geometric Ratios and Empirical Observations}
\label{sec:geometric}
% ============================================================

\subsection{The Reciprocal Silver Ratio}
Consider an object of height \(H\) casting a horizontal shadow of length \(S\). For a solar elevation angle \(\alpha\), the geometry dictates \(S/H = \cot(\alpha)\). Defining the angle from the vertical as \(\theta = 90^\circ - \alpha\), we obtain \(S/H = \tan(\theta)\). At \(\theta = 22.5^\circ\), the half-angle formula yields the exact identity:
\begin{equation}
\tan(22.5^\circ) = \sqrt{2} - 1 \approx 0.41421356.
\end{equation}

\subsection{Exploratory Field Measurements}
Five mango trees in Sana'a, Yemen, were measured for total height \(H\) and lower crown height \(h\). The vertical crown difference is \(C = H - h\), and the ratio is \(R = C/H\). 

\begin{table}[H]
\centering
\caption{Exploratory field measurements of tree crown ratios.}
\label{tab:tree_data}
\begin{tabular}{cccccc}
\toprule
Tree & \(H\) (m) & \(h\) (m) & \(C = H-h\) (m) & \(R = C/H\) & Deviation from \(\sqrt{2}-1\) \\
\midrule
1 & 7.80 & 4.50 & 3.30 & 0.42308 & \(+0.00887\) \\
2 & 8.50 & 5.00 & 3.50 & 0.41176 & \(-0.00245\) \\
3 & 7.20 & 4.20 & 3.00 & 0.41667 & \(+0.00245\) \\
4 & 9.00 & 5.40 & 3.60 & 0.40000 & \(-0.01421\) \\
5 & 7.50 & 4.40 & 3.10 & 0.41333 & \(-0.00088\) \\
\midrule
\multicolumn{4}{r}{\textbf{Mean \(R\):}} & \textbf{0.41297} & \\
\bottomrule
\end{tabular}
\end{table}

\begin{remark}
The mean \(R \approx 0.41297\) is numerically close to \(\sqrt{2}-1\). However, \(C\) represents a vertical height difference, not a horizontal shadow length. This table documents an empirical observation. It does not test the geometric law directly, and the numerical proximity does not prove the existence of a shared geometric law governing tree crowns. No causal link to the spectral operator is claimed.
\end{remark}

% ============================================================
\section{Algebraic Symmetry-Breaking}
\label{sec:algebraic}
% ============================================================

\subsection{The Kernel and the Gauging Identity}
Let \(u \in (0, 2\pi)\) and define the kernel:
\begin{equation}
\label{eq:g}
g(u) := \left(\frac{u}{2\pi}\right)^{5} \ln\!\left(\frac{2\pi}{u}\right) e^{-2\pi/(x_p u)}.
\end{equation}
Let \(p(u) := g'(u)/g(u)\). We define the first-order-corrected operator:
\begin{equation}
H_g := -\partial_u^2 + 2p(u)\partial_u + p'(u) - p(u)^2.
\end{equation}

\begin{proposition}[Gauging Identity]
For any \(Y \in C^2(0, 2\pi)\), \(H_g(gY) = -gY''\).
\end{proposition}
\begin{proof}
We expand the operator \(H_g\) directly acting on \(gY\):
\begin{align*}
H_g(gY) &= -(gY)'' + 2p(gY)' + (p'-p^2)gY \\
&= -gY'' - 2g'Y' - g''Y + 2p(gY' + g'Y) + (p'-p^2)gY \\
&= -gY'' - 2g'Y' - g''Y + 2\frac{g'}{g}(gY' + g'Y) + \left(\frac{g''}{g} - 2\frac{(g')^2}{g^2}\right)gY.
\end{align*}
Expanding the remaining terms:
\begin{align*}
H_g(gY) &= -gY'' - 2g'Y' - g''Y + 2g'Y' + 2\frac{(g')^2}{g}Y + g''Y - 2\frac{(g')^2}{g}Y \\
&= -gY''.
\end{align*}
The terms \(-2g'Y'\) and \(+2g'Y'\) cancel identically. The terms \(-g''Y\) and \(+g''Y\) cancel identically. The terms \(2(g')^2/g Y\) and \(-2(g')^2/g Y\) cancel identically, leaving only \(-gY''\).
\end{proof}

\subsection{Removal of Symmetry-Induced Cancellation}
In a symmetric model, the reflection \(\hat{P}u = 2\pi-u\) and the parity of states \(Y_n(u) = \sin(nu/2)\) force the transition matrix element \(\mathcal{A}_{13}^{(0)}\) (with \(D(u)=\cos(u/2)\)) to vanish.

However, \(g(u)\) is asymmetric. We define the unweighted \(L^2\) normalized states \(\Psi_n(u) = g(u)Y_n(u)/\sqrt{N_n}\), where \(N_n = \int_0^{2\pi} g(u)^2 Y_n(u)^2 du\). We compute:
\begin{equation}
\mathcal{A}_{mn} = \int_0^{2\pi} \Psi_m D \Psi_n \, du, \qquad R_{13/12} = \frac{|\mathcal{A}_{13}|^2}{|\mathcal{A}_{12}|^2}.
\end{equation}
Using numerical integration (Simpson quadrature and adaptive methods) at \(x_p=21\), the asymmetry breaks the cancellation, yielding a computed numerical value (independently reproduced using the supplied Python code):
\begin{equation}
R_{13/12} \approx 0.6035155248.
\end{equation}

% ============================================================
\section{Towards a Spectral Problem}
\label{sec:spectral}
% ============================================================

\subsection{The Formal Zero Mode and Factorization}
The gauging identity motivates the definition of a Schr\"odinger-like operator:
\begin{equation}
L := -\partial_u^2 + V_{\text{eff}}(u), \qquad V_{\text{eff}}(u) := \frac{g''(u)}{g(u)}.
\end{equation}
\begin{proposition}[Formal Zero Mode]
The kernel \(g\) satisfies \(Lg = 0\). Furthermore, \(L\) admits a supersymmetric (SUSY QM) factorization:
\begin{equation}
L = (-\partial_u - p)(\partial_u - p).
\end{equation}
\end{proposition}

\begin{remark}
The identity \(Lg=0\) is a formal algebraic statement. It does not assert that \(g\) is a normalizable eigenfunction on a specified self-adjoint domain. Furthermore, the algebraic factorization alone does not guarantee the positivity of the operator or dictate the spectrum; those properties depend critically on the choice of boundary conditions.
\end{remark}

\subsection{Asymptotics of the Effective Potential}
The spectral properties are constrained by the singularities of \(V_{\text{eff}}\):
\begin{itemize}
    \item \textbf{Near \(u \to 0^{+}\):} \(p(u) \sim 2\pi/(x_p u^2)\), so \(V_{\text{eff}}(u) \sim \frac{4\pi^2}{x_p^2 u^4} \to +\infty\).
    \item \textbf{Near \(u \to 2\pi^{-}\):} Let \(t = 2\pi - u \to 0^+\). Expanding \(g(u) = At + Bt^2 + O(t^3)\), we find \(A = \frac{e^{-1/x_p}}{2\pi}\) and \(B = -\frac{e^{-1/x_p}}{8\pi^2}(9 + \frac{2}{x_p})\). Since \(g''/g \sim 2B/(At)\), we obtain:
    \begin{equation}
    V_{\text{eff}}(u) \sim -\frac{9 + 2/x_p}{2\pi(2\pi - u)}.
    \end{equation}
    For \(x_p = 21\), this yields \(V_{\text{eff}}(u) \sim -\frac{191}{42\pi(2\pi - u)} \to -\infty\).
\end{itemize}
These asymmetric singularities dictate the domain and boundary conditions required for self-adjoint extensions of \(L\), which remain an open question.

% ============================================================
\section{Summary and Remaining Work}
\label{sec:summary}
% ============================================================

To ensure the integrity of this exploratory work, we explicitly demarcate what has been accomplished, what requires further verification, and what remains open. 

\paragraph{What has been accomplished.}
We have presented the mathematical definitions and formulations of the four frameworks. We established the analytical gauging identity and the formal zero-mode factorization. We provided a numerical evaluation of the transition ratio \(R_{13/12}\) and documented the exploratory geometric measurements. 

\paragraph{What requires verification.}
A rigorous line-by-line audit of the traceability matrix (Appendix~\ref{app:audit}) against the original HTML frames is necessary. Furthermore, all bibliographic references and DOI links must be verified prior to final publication.

\paragraph{What remains open.}
The complete spectral interpretation of the operator \(L\) remains unresolved, specifically concerning the self-adjoint domain, boundary conditions, and the normalizability of the kernel \(g(u)\). Additionally, the empirical tree measurements require further refinement to serve as a direct test of the geometric shadow law.

\paragraph{What is not claimed.}
It is not claimed that these four parallel paths have established a unified physical law, a causal link between the distinct frames, or a proof of the Riemann Hypothesis. The numerical coincidence observed in the tree crown ratios does not constitute a proof of a natural law.

% ============================================================
\section*{Epilogue: A Message to the Next Researcher}
\label{sec:epilogue}
% ============================================================

\textit{This document serves to preserve a starting point, not to serve as a final verdict on the validity of every inference. Whoever continues this project must strictly separate mathematical proof, computational results, and empirical observation; re-verify every testable result; and never attribute more to the project than the evidence allows. The value of this work lies in the ideas it makes possible to examine, and in the questions it leaves clear for those who follow.}

% ============================================================
\appendix
\section{Traceability Matrix}
\label{app:audit}
% ============================================================

\begin{table}[H]
\centering
\caption{Preliminary Traceability Matrix: Mapping of results from original frames to the consolidated draft.}
\label{tab:traceability}
\scriptsize
\begin{tabular}{p{2.8cm} p{2.2cm} p{2.5cm} p{6cm}}
\toprule
\textbf{Result / Claim} & \textbf{Source Frame} & \textbf{Action in Draft} & \textbf{Verification Status \& Notes} \\
\midrule
Sign analysis of \(Z(x)\) & Frame 1 & Re-derived & \textbf{Proven.} Rigorous mathematical check on \((1,2)\) and \((2,\infty)\). \\
\midrule
Global Minimum at \(x=2\) & Frame 1 & Re-derived & \textbf{Proven.} Follows directly from \(A'(x)=Z(x)\). (Note: \(x=42\) max was discarded). \\
\midrule
Exact ratio \(\tan(22.5^\circ) = \sqrt{2}-1\) & Frame 2 & Verbatim & \textbf{Proven.} Standard trigonometric identity. \\
\midrule
Tree crown ratio \(R \approx 0.41297\) & Frame 2 & Verbatim & \textbf{Observed.} \(N=5\) sample. Vertical heights only; not a shadow test. \\
\midrule
Gauging Identity \(H_g(gY) = -gY''\) & Frame 3 & Re-derived & \textbf{Proven.} Explicit algebraic expansion provided in proof. \\
\midrule
Numerical ratio \(R_{13/12} \approx 0.60352\) & Frame 3 & Re-computed & \textbf{Computed \& Verified.} The supplied Python code was executed successfully, and its results matched the reported values. \\
\midrule
Zero-mode \(Lg = 0\) & Frame 4 & Re-derived & \textbf{Proven (Formal).} Does not imply normalizability. \\
\midrule
Factorization of \(L\) & Frame 4 & Re-derived & \textbf{Proven (Algebraic).} Does not prove operator positivity. \\
\midrule
Asymptotic \(V_{\text{eff}} \sim -\frac{191}{42\pi t}\) & Frame 4 & Re-derived & \textbf{Proven.} Explicit Taylor expansion provided. \\
\bottomrule
\end{tabular}
\end{table}

\section{Reproducibility: Python Code}
\label{app:code}
% ============================================================

The following Python code reproduces the numerical results in Section~\ref{sec:algebraic}. It has been independently executed in a standard Python environment with \texttt{NumPy} and \texttt{SciPy}, yielding the expected output.

\begin{verbatim}
import sys
import numpy as np
from scipy.integrate import quad

XP = 21.0
PI2 = 2.0 * np.pi

EXPECTED_A12 = 0.6220748484804
EXPECTED_A13 = -0.4832666943981
EXPECTED_R = 0.6035155247873

EXPECTED_TABLE = {
    20: 0.6033930653531,
    40: 0.6035132002961,
    80: 0.6035154874047,
    200: 0.6035155246331,
    400: 0.6035155247849,
}

def g(u, xp=XP):
    if u <= 0.0 or u >= PI2:
        return 0.0
    return ((u / PI2)**5 * np.log(PI2 / u) * np.exp(-PI2 / (xp * u)))

def Y(n, u): return np.sin(n * u / 2.0)
def D(u): return np.cos(u / 2.0)

def simpson(f, a, b, N):
    if N < 2: raise ValueError("N must be at least 2")
    if N % 2: N += 1
    h = (b - a) / N
    total = f(a) + f(b)
    total += 4.0 * sum(f(a + i * h) for i in range(1, N, 2))
    total += 2.0 * sum(f(a + i * h) for i in range(2, N, 2))
    return total * h / 3.0

def matrix_elements(integrate, xp=XP):
    norms = [integrate(lambda u: (g(u, xp) * Y(n, u))**2) for n in (1, 2, 3)]
    if any(not np.isfinite(v) or v <= 0 for v in norms):
        raise ArithmeticError("Invalid normalization integral")

    def amplitude(m, n):
        return integrate(lambda u: (g(u, xp) * Y(m, u) / np.sqrt(norms[m - 1])) * 
                         D(u) * (g(u, xp) * Y(n, u) / np.sqrt(norms[n - 1])))
    a12 = amplitude(1, 2)
    a13 = amplitude(1, 3)
    if not np.isfinite(a12) or not np.isfinite(a13) or a12 == 0:
        raise ArithmeticError("Invalid matrix element")
    return a12, a13, (a13 / a12)**2

def adaptive_integral(f):
    value, error = quad(f, 0.0, PI2, limit=300, epsabs=1e-12, epsrel=1e-12)
    if not np.isfinite(value) or not np.isfinite(error):
        raise ArithmeticError("Adaptive quadrature failed")
    return value

def simpson_integral(f, N): return simpson(f, 0.0, PI2, N)

def main():
    a12, a13, r_ref = matrix_elements(adaptive_integral)
    assert abs(a12 - EXPECTED_A12) < 5e-11
    assert abs(a13 - EXPECTED_A13) < 5e-11
    assert abs(r_ref - EXPECTED_R) < 5e-11

    print(f"Adaptive quad: A12={a12:.13f}, A13={a13:.13f}, R={r_ref:.13f}")
    
    previous_error = float("inf")
    for N, expected in EXPECTED_TABLE.items():
        _, _, r = matrix_elements(lambda f, n=N: simpson_integral(f, n))
        error = abs(r - r_ref)
        assert abs(r - expected) < 5e-11
        assert error < previous_error
        print(f"N={N:3d} | R={r:.13f} | Error={error:.4e}")
        previous_error = error
        
    print("\nVerification successful.")

if __name__ == "__main__":
    try:
        main()
    except Exception as exc:
        print(f"NOT VERIFIED: {exc}", file=sys.stderr)
        sys.exit(1)
\end{verbatim}

% ============================================================
\begin{thebibliography}{99}

\bibitem{Frame1}
A.~M.~N.~Al-Jabri, \emph{The Numerical Framework: A Damping Model, the Riemann Zeta Zeros, and the Proposed Integer Cutoff $x_p = 21$}, Al-Hindara Working Paper No.~1, Zenodo (2026). DOI: \href{https://doi.org/10.5281/zenodo.23540283}{10.5281/zenodo.23540283}.

\bibitem{Frame2}
A.~M.~N.~Al-Jabri, \emph{The Tree and the Shadow: Geometric Derivation of the Reciprocal Silver Ratio and an Exploratory Statistical Analysis of Tree Crown Ratios}, Al-Hindara Working Paper No.~2, Zenodo (2026). DOI: \href{https://doi.org/10.5281/zenodo.23540796}{10.5281/zenodo.23540796}.

\bibitem{Frame3}
A.~M.~N.~Al-Jabri, \emph{Frame 3: Algebraic Transformation and Symmetry-Induced Cancellation Removal}, Al-Hindara Project Technical Documentation (2026).

\bibitem{Frame4}
A.~M.~N.~Al-Jabri, \emph{Frame 4: Towards a Spectral Problem, Zero Mode, and Zeta Regularization}, Al-Hindara Project Technical Documentation (2026).

\bibitem{Teschl}
G.~Teschl, \emph{Mathematical Methods in Quantum Mechanics: With Applications to Schr\"odinger Operators}, 2nd ed., American Mathematical Society, Providence, RI, 2014.

\end{thebibliography}

\end{document}